\section{Related Work}
\label{sec:related}

Research on autism spectrum disorder (ASD) classification using resting-state functional magnetic resonance imaging (rs-fMRI) has gradually developed from statistical modelling and functional-connectivity representation. More recent research is focused on multisite learning, deep architectures, dynamic connectivity, multimodal fusion, and reliability-aware prediction. In early studies, methods are established for representing high-dimensional covariance and brain-network information. These methods are later extended through deep learning, multiple brain atlases, temporal modelling, and graph-based approaches to improve classification performance. Alongside these developments, calibration, interpretability, selective prediction, and statistical validation are increasingly used to improve the reliability of machine-learning predictions. Overall, the performance and reliability of ASD classification are shaped by several methodological choices, including how functional connectivity is represented, how variation across acquisition sites is controlled, and how prediction uncertainty is assessed.

\subsection{Foundations of Functional Connectome-Based Analysis}


The methodological basis for connectome-based ASD classification is established through a sequence of statistical and neuroimaging developments. Reliable evaluation first requires uncertainty in classification error to be quantified, for which exact confidence limits for binomial outcomes are provided. Reliable connectivity estimation is then required before such models can be evaluated. High-dimensional covariance estimation is stabilised through the Ledoit--Wolf shrinkage method \citep{LedoitWolf2004}, which is particularly relevant to rs-fMRI because many regional signals are estimated from limited time points. Once stable features are obtained, reliable comparison between classifiers is supported by statistical procedures for model evaluation. The use of multisite data introduces an additional source of variation, which is addressed through the ComBat empirical-Bayes framework for batch-effect correction \citep{Johnsonetal2007} and its later adaptation to neuroimaging for scanner- and site-related harmonisation \citep{Fortinetal2018}. With these statistical foundations in place, the structure of connectivity itself is modelled more directly. Covariance matrices are represented according to their non-Euclidean geometry \citep{Varoquauxetal2010}, from which tangent-space representations are established for functional-connectivity analysis. Biological interpretation is then supported by the identification of seven large-scale intrinsic brain networks \citep{Yeoetal2011}, through which individual connections are interpreted within broader functional systems. Connectivity is first represented statically across the full scan, after which temporal variation is examined through dynamic functional connectivity \citep{Hutchisonetal2013}. Functional connectivity is therefore developed from stable covariance estimation toward geometrically informed, biologically organised, and temporally resolved representations.

\subsection{Multisite ASD Data and Reproducibility}

Large-scale ASD neuroimaging research is substantially advanced through the release of ABIDE \citep{DiMartinoetal2014}. Data from multiple international sites are combined, making larger ASD classification studies possible. This increased scale also introduces variation in scanner type, acquisition protocol, demographic composition, and site characteristics. As a result, prediction models risk capturing acquisition-related differences instead of ASD-related neural patterns. The use of larger multisite datasets also requires high-dimensional brain information to be represented carefully. Regularised matrix regression is proposed for this purpose, where predictors are kept in matrix form so that relationships among features are preserved during modelling \citep{ZhouLi2014}. Careful analysis is also required when changes in connectivity over time are examined. Sliding-window estimates are shown to depend on window length, and apparent connectivity changes are potentially produced by analysis choices \citep{LeonardiVanDeVille2015}. Temporal information therefore needs to be evaluated with careful control of scan length and window design. These methodological choices directly affect whether findings are reproducible across datasets and sites. Using multisite resting-state data, biomarker performance is shown to depend on preprocessing, feature construction, and validation design \citep{Abrahametal2017}. The reliability of time-varying connectivity is also examined through a review of dynamic functional-connectivity methods \citep{Pretietal2017}. Together, these studies show that stable ASD-related findings depend on careful control of participants, analysis choices, and acquisition sites.

\subsection{Reliability, Calibration, and Explainability}

Alongside developments in neuroimaging representation and generalisation, reliability-oriented machine-learning methods are also developed. Probability calibration is examined in a study \citep{Guoetal2017}, where high classification accuracy is shown to be insufficient for reliable probability estimates. Participants may be ranked correctly even when the predicted probabilities do not match observed outcome frequencies. Calibration is therefore assessed through calibration-in-the-large, calibration slope, and the Brier score. Systematic probability offsets can be corrected through recalibration fitted on data separate from the evaluation set \citep{VanCalsteretal2019,Steyerbergetal2010,ZadroznyElkan2002}. Selective classification is introduced in a study \citep{GeifmanElYaniv2017}, where uncertain cases are allowed to remain unclassified instead of requiring a prediction for every input. Prediction error can therefore be reduced at the cost of lower coverage, which is relevant to screening settings where uncertain cases are referred for further clinical assessment. Model interpretation is also considered at both the individual-prediction and feature-pattern levels. Contributions to individual predictions are quantified through SHAP \citep{LundbergLee2017}. For linear models, model weights are not interpreted directly as differences between groups, and a corresponding pattern relating each feature to the model output is required \citep{Haufeetal2014}. Evaluation is consequently extended beyond discrimination: calibration is used to assess probability reliability, selective prediction is used to define when a prediction is issued, and explanation methods are used to identify the information contributing to model output. Clinical usefulness is assessed separately through decision-curve analysis \citep{VanCalsteretal2018}. Reliable ASD classification is therefore evaluated through discrimination, calibration, selective decision-making, interpretability, and clinical usefulness.



\subsection{Deep Learning for ASD Classification}

Alongside these reliability-focused developments, deep learning is increasingly applied to ABIDE functional-connectivity data. Distributed connectivity patterns are learned automatically and an accuracy of approximately 70.0\% is reported, showing that ASD-related information is captured through deep feature learning \citep{Heinsfeldetal2018}. The effect of multisite variation still needs to be addressed when such models are applied to ABIDE. For this purpose, ComBat is adapted for neuroimaging harmonisation in a study \citep{Fortinetal2018}, allowing scanner- and site-related variation to be reduced while biological variation is preserved. The harmonisation step also needs to be estimated without information from held-out test data. Deep representation learning is further developed through ASD-DiagNet, which is proposed in a study \citep{Eslamietal2019}. Representation learning and classification are combined, and an accuracy of approximately 70.3\% is reported. The small difference in performance indicates that classification is also potentially influenced by the way functional connectivity is represented. This role of connectivity representation is examined more directly in a controlled benchmarking study \citep{Dadietal2019}. Multiple connectome representations and predictive models are compared under the same setting, and tangent-space connectivity is identified as a strong representation. With an appropriate connectivity representation, comparatively simple classifiers are shown to remain competitive with more complex nonlinear models. Sparse representation learning is further explored through ASD-SAENet, which is proposed in a study \citep{AlmuqhimSaeed2021}. A sparse autoencoder is combined with a deep neural network, and an accuracy of approximately 70.8\% is reported. Across these deep-learning approaches, the increase from 70.0\% to 70.3\% and then 70.8\% remains relatively small. This pattern suggests that deeper representation learning alone does not necessarily fully address the challenges of heterogeneous multisite rs-fMRI data.


\subsection{Evidence from Systematic Reviews}

The broader performance of rs-fMRI machine-learning methods is examined through a systematic review and meta-analysis in \cite{Santanaetal2022}. Moderate pooled sensitivity, specificity, and discrimination are reported, while considerable methodological heterogeneity is identified. Differences in preprocessing, sample selection, validation strategy, feature extraction, and classifier design make direct comparison between studies difficult. A later systematic review of deep-learning and rs-fMRI biomarkers \cite{Hudaetal2024} reports continued methodological development but identifies similar limitations related to generalisability, reproducibility, and inconsistent evaluation. These reviews provide an important analytical perspective: higher accuracy reported by an individual model cannot be interpreted independently of the dataset definition and experimental pipeline used to obtain it. Therefore, later ASD classification studies increasingly attempt to improve the representation itself rather than only increasing classifier depth.

\subsection{Multi-Atlas and Dynamic Connectivity Models}

Dependence on a single brain parcellation is addressed through MADE-for-ASD \cite{Liuetal2024MADE}. A multi-atlas deep ensemble is introduced, and approximately 75.2\% accuracy is reported. Different atlases describe the same brain according to different regional boundaries, so their predictions potentially contain complementary information. Multi-atlas fusion therefore provides a possible way to reduce dependence on one predefined parcellation. However, the improvement is introduced together with a more complex ensemble architecture. It therefore becomes difficult to determine whether performance is improved because of atlas diversity, deep-model capacity, demographic information, or their interaction. This motivates controlled evaluations in which atlas fusion is varied while the remaining pipeline is kept unchanged. Dynamic functional connectivity is later revisited through a GNN-LSTM framework \cite{Tangetal2025GNNLSTM}. Graph neural networks are combined with temporal sequence modelling, and approximately 80.4\% accuracy is reported on ABIDE-I. Compared with static approaches, connectivity variation across the scan is explicitly represented rather than being averaged into one matrix. This development addresses an important limitation of static functional connectivity, but it introduces several additional decisions. Window length, overlap, temporal ordering, graph construction, and sequence-model complexity all affect performance. Therefore, the contribution of temporal variation requires separation from the contribution of the classifier through controlled ablation.

\subsection{Explainable and Multimodal ASD Models}

Interpretability receives greater attention in recent ASD classification research. Explainable artificial intelligence is applied to ABIDE fMRI data in a study \cite{Vidyaetal2025}, where highly predictive performance is reported under the stated filtering condition and influential brain regions are identified through multiple explanation methods. Such work demonstrates the potential of connecting prediction with biologically interpretable patterns. However, unusually high performance on heterogeneous multisite data also increases the importance of strict validation. Feature selection, filtering, preprocessing, and cross-validation are required to remain independent of held-out data. The broader lesson from earlier benchmarking \cite{Dadietal2019} and meta-analysis \cite{Santanaetal2022} therefore remains relevant: reported performance requires interpretation together with the complete experimental procedure. The progression toward richer representations continues through CAS-GNN \cite{Tangetal2026CASGNN}. Structural and functional MRI are integrated using cross-attention and subject-adaptive graph construction, and approximately 79.3\% accuracy is reported. Compared with earlier functional-connectivity-only approaches, multimodal information is used to represent each participant more comprehensively. Although multimodal graph learning increases representational capacity, additional learned components also make generalisation more difficult to assess. Performance gains therefore require evaluation against cross-site stability, model complexity, and potential overfitting.

\subsection{Research Gap}

Two connected methodological developments are identified in the literature. In the first, functional connectivity is progressively represented through stable covariance estimation \cite{LedoitWolf2004}, geometric modelling \cite{Varoquauxetal2010}, network-level organisation \cite{Yeoetal2011}, and dynamic connectivity \cite{Hutchisonetal2013,LeonardiVanDeVille2015,Pretietal2017}. These representations are extended to multisite analysis \cite{DiMartinoetal2014}, reproducibility assessment \cite{Abrahametal2017}, deep learning \cite{Heinsfeldetal2018,Eslamietal2019,AlmuqhimSaeed2021}, multi-atlas fusion \cite{Liuetal2024MADE}, dynamic graph modelling \cite{Tangetal2025GNNLSTM}, and multimodal graph learning \cite{Tangetal2026CASGNN}. Increasingly rich representations are therefore introduced, together with additional sources of variation and overfitting that require controlled evaluation. In the second development, greater attention is given to prediction reliability.Confidence estimation and statistical comparison are formalised and reliability is further addressed through harmonisation \cite{Johnsonetal2007,Fortinetal2018}, calibration \cite{Guoetal2017}, selective prediction \cite{GeifmanElYaniv2017}, explainability \cite{LundbergLee2017}, and decision-oriented evaluation \cite{VanCalsteretal2018}. Classification accuracy alone is therefore insufficient for screening-oriented use; reliable probabilities, uncertainty-aware decisions, interpretable evidence, and performance across unseen sites are also required. These methodological directions are commonly evaluated separately, leaving dynamic connectivity, multi-atlas information, harmonisation, interpretability, calibration, selective prediction, and cross-site generalisation without a common controlled evaluation. The benefit of increasingly complex classifiers over carefully designed functional-connectivity representations also remains uncertain when identical validation conditions are imposed. The present work is therefore structured as a controlled, reliability-oriented evaluation. Complementary functional-connectivity representations are compared using the same data partitions, tangent-space quantities are estimated from training data only, and cross-site performance is assessed through leave-one-site-out validation. Ensemble probabilities are assessed for calibration and are recalibrated within the training partitions when systematic offset is identified. The ensemble score is also attributed exactly to functional-system pairs, and the reproducibility of this attribution is tested against a permutation reference. Each component is evaluated independently so that observed improvements are attributed to the corresponding methodological choice rather than to model complexity alone.

